\subsubsection{検証・妥当性確認・認証}

設計の分析と評価は、検証、妥当性確認、認証に関わる。

\par\smallskip\noindent\refstepcounter{table}\label{tab:ch03-07}\textbf{表 \thetable: 第03章 ソフトウェア設計 表07}\par\smallskip
\begin{tabularx}{\linewidth}{L{32mm}Y}
\toprule
\textbf{活動} & \textbf{意味} \\
\midrule
検証 & 設計が、明示された要求を満たしているかを確認することである。 \\
妥当性確認 & その設計で作られるシステムが、顧客、利用者、運用者、保守担当者の期待を満たせるかを確認することである。 \\
認証 & 第三者が、設計が仕様や意図された利用に適合していると証明または表明することである。 \\
\bottomrule
\end{tabularx}

検証は「要求どおりか」を問う。妥当性確認は「本当に役に立つか」を問う。認証は「決められた基準に適合していると第三者が認めるか」を問う。安全性や規制が重要な領域では、設計段階から認証を見据えた証拠を残す必要がある。
